\section{Verification note}\label{sec:verification-note}

The formal results checked in Lean 4 are the estimator guarantee in \cref{obj:thm:rectangular-upper}, the supplied-propensity benchmark in \cref{obj:thm:supplied-propensity}, the smooth-frame construction certificate in \cref{obj:thm:marked-component-certificate}, and the nuisance lower bound in \cref{obj:lem:nuisance-frontier-converse}. Their combination yields the minimax sandwich in \cref{obj:thm:sharp-annotation-frontier} and the acquisition consequences in \cref{obj:prop:annotation-threshold}. Auxiliary checked results cover local polynomial projection, population positivity, rectangular sampling error, and Hellinger calculus in \cref{obj:lem:local-polynomial-projection-facts,obj:lem:population-positivity,obj:lem:rectangular-sampling-bound,obj:lem:hellinger-block-calculus}.

The companion artifact for this paper version, whose working-paper number and version appear in the first-page stamp, is available from \url{https://causalsmith.org/papers/stat_twosample_pointcate_annotation_frontier_v1}. Its interactive formal links provide the theorem-to-file crosswalk and exact source excerpts. In the source archive, the checked modules are rooted at \texttt{CausalSmith/Stat/STAT\_TwosamplePointcateAnnotationFrontier\_Research}; the generated file \texttt{presentation\_crosswalk.json} records the mapping from every displayed result to its Lean module and declaration. The crosswalk is regenerated with each paper version, so it describes the results as printed in this version.

The primitive-model and sampling restrictions are assumed inputs. These comprise the known uniform covariate design, conditional exchangeability, a fixed public overlap level \(\varepsilon\in(0,1/2)\), binary potential responses, H\"older restrictions, and probability-valued arm-mean bounds defining \cref{obj:def:primitive-class}, together with the independent product experiment and randomizer in \cref{obj:ass:sample-law}. Each guarantee retains its stated public-parameter, sample-size, scale, amplitude, and occupancy conditions.

The absolute-loss testing result in \cref{obj:lem:published-fuzzy-testing} is proved in the formal artifact without external formal assumptions; \citet[Section 3, Lemma 1]{KennedyBalakrishnanRobinsWasserman2024} remains its bibliographic source and statistical motivation. For the supervised comparison, \cref{obj:lem:published-cate-benchmark} checks the numerical exponent identities, smoothness cutoff, and piecewise-rate equality associated with \citet[Section 3, Theorem 1]{KennedyBalakrishnanRobinsWasserman2024}. The published statistical lower bound belongs to that source's model and assumptions. The minimax interpretation for the present known-uniform primitive class is established by \cref{obj:thm:sharp-annotation-frontier}, through this paper's estimator guarantee and original-record lower bounds. The formal numerical identity also assigns totalized negative powers the value zero at \(n=0\); every statistical risk and acquisition guarantee in the paper uses \(n\ge2\), so this representation value has no precision interpretation.

For the machine-readable crosswalk, the verified rate and decision are also packaged in the following representation-level handle; no statistical argument depends on this packaging.

\begin{definitionv}{def:frontier-handle}[Rate and decision handle]
The frontier handle is the ordered pair \((r_*,T_*)\), parameterized by \(d,n,m\in\mathbb N\) and \(\alpha,\beta,\gamma,\varepsilon\in\mathbb R\), where \(\varepsilon\) is the overlap parameter. Its components are the sharp rate and decision of \cref{obj:def:sharp-rate,obj:def:sharp-decision}:
\[
r_*=
\max\left\{
n^{-\gamma/(2\gamma+d)},
\bigl(n(n+m)\bigr)^{-\gamma/(2\gamma+d+\gamma d/(\alpha+\beta))}
\right\},
\qquad
T_*(\mathcal D_{n,m},U)
=
\widehat T_{\mathrm{up}}(\mathcal D_{n,m}).
\]
Here \(\mathcal D_{n,m}\) is the original dataset and \(U\in\mathbb R\) is the randomizer coordinate. The decision component uses the publicly tuned estimator of \cref{obj:def:tuning} with parameters \(d,\alpha,\beta,\gamma,\varepsilon,n,m\).
\end{definitionv}
